\documentclass[11pt]{article}
\usepackage[margin=1.05in]{geometry}
\usepackage{amsmath, amssymb, booktabs}
\usepackage[colorlinks=true, linkcolor=blue, citecolor=blue]{hyperref}

\title{The computational tasks behind the layered-memory codes\\
\large what is computed, by which algorithm, at what cost,
and where it could be much faster}
\author{}
\date{July 2026 --- working notes}

\newcommand{\prof}{\mathbf{m}}
\newcommand{\field}[1]{\paragraph{#1.}}
\newcommand{\Exp}{\mathrm{Exp}}

\begin{document}
\maketitle

\noindent
Every experiment in this project is assembled from eight
computational tasks.  This note describes each task in a fixed
format: what the task is, the formula it evaluates, its inputs and
outputs, the algorithm we currently use, the cost of that algorithm
and how the cost grows, and ideas for computing the same thing
faster.  Ideas that give identical results are marked [exact]; ideas
that introduce a controlled error are marked [approx].
Section~\ref{sec:defs} defines all concepts and symbols; nothing is
used before it is defined there.

\section{Definitions}
\label{sec:defs}

\field{The data}
The data is a sequence $x_1, x_2, \dots, x_n$ of \emph{symbols}.  In
our experiments a symbol is an English word, and the sequence is the
corpus \texttt{text8} with $n = 1.7 \cdot 10^7$ words.  There are $d$
possible symbols (the \emph{alphabet size}); depending on the
experiment, $d$ ranges from $256$ to the full number of distinct
words, $253{,}855$.

\field{States}
A predictor with memory works as follows.  A fixed rule $\sigma$
assigns to every position $t$ a \emph{state}
$s_t = \sigma(x_1, \dots, x_{t-1})$, computed from the symbols before
position $t$.  Examples of rules used in this project: the state is
the symbol immediately before position $t$; or the pair of the two
preceding symbols; or the symbol that occurred $\delta$ positions
earlier (we call $\delta$ the \emph{lag}).  The predictor keeps one
table of statistics per state, and at position $t$ it predicts $x_t$
using the table that belongs to the state $s_t$.

\field{Count rows}
Fix a state $s$.  Among the positions $t$ whose state equals $s$,
count, for each symbol $x$, how many of these positions actually
carried the symbol $x$:
\[
  m_x(s) \;=\; \#\{\, t \;:\; s_t = s \text{ and } x_t = x \,\}.
\]
The vector $\bigl(m_1(s), \dots, m_d(s)\bigr)$ is called the
\emph{row} of the state $s$.  In words: the row of $s$ records how
often each symbol occurred at the positions where the state was $s$.
We write $s$ for the number of nonzero entries of a row when no
confusion with the state is possible, $N = \sum_x m_x$ for the sum of
the row's entries, and we call a symbol \emph{observed in the row} if
its entry is nonzero.

\field{Profiles}
The \emph{profile} of a row is the list of its nonzero entries,
sorted in increasing order, with the identities of the symbols
discarded.  For example, the row $(0, 3, 0, 1, 1)$ has the profile
$(1, 1, 3)$.  We write a profile as $\prof = (m_1, \dots, m_k)$ with
$0 < m_1 \le \dots \le m_k$.  Two further quantities of a profile
matter for costs: $\tilde k$, the number of \emph{distinct} values
among its entries (the example has $\tilde k = 2$: the values $1$
and $3$), and the multiplicity $\mu_c$ of each distinct value $c$
(here $\mu_1 = 2$, $\mu_3 = 1$).

\field{The estimator}
All experiments use one estimator of probabilities from counts, the
\emph{layered mixture} of the companion paper.  It is a Bayesian
mixture: a prior distribution over probability vectors, updated by
the observed counts.  The prior is built in \emph{levels}.  At level
$L$, attach to each symbol $x$ a random weight
\[
  W_x \;=\; E_{x,1} \cdot E_{x,2} \cdots E_{x,L},
\]
where all $E_{x,\ell}$ are independent random variables with the
standard exponential distribution (density $e^{-v}$ for $v > 0$).
The random probability vector of level $L$ is
$p_x = W_x / (W_1 + \dots + W_d)$.  The larger $L$ is, the more
uneven these probability vectors tend to be, so $L$ acts as a
roughness parameter.  The full estimator does not choose $L$; it
averages uniformly over $L = 1, \dots, L_{\max}$, with
$L_{\max} = \lceil 2 c^{*} \ln d \rceil$ and $c^{*} \approx 2.365$ a
constant from the companion paper.  In our runs $L_{\max}$ is between
$33$ and $59$.

\field{The central quantity}
For a profile $\prof$ and a level $L$, define
\begin{equation}
  q_L(\prof) \;=\; \mathbb{E}\Bigl[\, p_{y_1}^{\,m_1} \cdot
  p_{y_2}^{\,m_2} \cdots p_{y_k}^{\,m_k} \Bigr],
  \qquad
  q_{\mathrm{avg}}(\prof) \;=\; \frac{1}{L_{\max}}
  \sum_{L=1}^{L_{\max}} q_L(\prof),
  \label{eq:qdef}
\end{equation}
where $y_1, \dots, y_k$ are any $k$ distinct symbols and the
expectation is over the random vector $p$ of level $L$.  Because the
prior treats all symbols in the same way (it is exchangeable), the
value does not depend on which distinct symbols are chosen --- only
the profile matters.  This is the number the whole project runs on:
$q_{\mathrm{avg}}(\prof)$ is the probability that the estimator
assigns to any one particular sequence of symbols whose counts form
the profile $\prof$.  The compressed length contributed by one row is
$-\log_2 q_{\mathrm{avg}}$ of its profile; the length of a whole
model is the sum of this over its rows.  Because rows with equal
profiles contribute equal amounts, all computations deduplicate rows
by profile first; this typically shrinks the workload by a factor of
$50$ to $100$.

\field{The integral form}
Eq.~\eqref{eq:qdef} is an expectation over $d$ random weights; to
compute it we turn it into a one-dimensional integral.  The
denominator $(W_1 + \dots + W_d)^N$ is removed with the identity
$A^{-N} = \frac{1}{\Gamma(N)} \int_0^\infty t^{N-1} e^{-tA}\, dt$
(valid for $A > 0$; here $t$ is a real integration variable, a
scalar).  After this substitution the expectation splits into a
product of one factor per symbol, because the weights of different
symbols are independent.  The result is exact:
\begin{equation}
  q_L(\prof) \;=\; \frac{1}{\Gamma(N)} \int_0^{\infty} t^{N-1}\,
  \varphi^{(L)}_0(t)^{\,d-k}\,
  \prod_{i=1}^{k} \varphi^{(L)}_{m_i}(t)\; dt,
  \label{eq:integral}
\end{equation}
where the function $\varphi^{(L)}_r$, which we call the
\emph{moment function}, is defined by
\begin{equation}
  \varphi^{(L)}_{r}(t) \;=\; \mathbb{E}\bigl[\,Y^{\,r}\, e^{-tY}
  \bigr],
  \qquad Y = E_1 E_2 \cdots E_L
  \label{eq:phidef}
\end{equation}
with $E_1, \dots, E_L$ independent standard exponentials.  In words:
$Y$ is a product of $L$ independent exponential random variables, and
$\varphi^{(L)}_r(t)$ is the expected value of $Y^r e^{-tY}$ --- a
function of the single real argument $t \ge 0$, one such function for
each pair $(L, r)$.  The factor $\varphi^{(L)}_0(t)^{d-k}$ in
Eq.~\eqref{eq:integral} collects the $d - k$ symbols that are not
observed in the row (their count is $r = 0$).

\field{The computational axis}
All numerics happen on the logarithmic axis $u = \ln t$ ($u$ is a
scalar), restricted to a finite grid $u_1 < u_2 < \dots < u_G$; in
the largest runs $G = 3.6 \cdot 10^4$.  Taking logarithms of the
integrand of Eq.~\eqref{eq:integral} and substituting $t = e^u$
gives the \emph{log-integrand}
\begin{equation}
  \psi_L(u) \;=\; N u \;+\; (d - k) \ln \varphi^{(L)}_0(e^u)
  \;+\; \sum_{\text{distinct values } c \text{ of } \prof}
  \mu_c \ln \varphi^{(L)}_c(e^u),
  \label{eq:psi}
\end{equation}
so that $q_L(\prof) = \frac{1}{\Gamma(N)} \int_{-\infty}^{\infty}
e^{\psi_L(u)}\, du$.  Two exact facts anchor the numerics.  First,
at level one the moment function has the closed form
$\varphi^{(1)}_r(t) = \Gamma(r+1) / (1+t)^{r+1}$, and the whole
integral at $L = 1$ reduces to a classical closed form (the
Dirichlet--multinomial); no numerics are needed there.  Second, as
$t \to 0$ the moment function tends to the constant
$\Gamma(r+1)^L$; therefore, to the left of the grid, $\psi_L(u) =
Nu + \text{const} + O(e^u)$, and the left tail of the integral is
available in closed form.

\bigskip
\noindent
Symbols used in the cost formulas, with the values measured in our
runs:

\begin{center}
\small
\begin{tabular}{llr}
\toprule
symbol & meaning & measured \\
\midrule
$n$ & length of the corpus, in symbols & $1.7\cdot 10^7$ \\
$d$ & alphabet size & $256$ -- $2.6\cdot 10^5$ \\
$L_{\max}$ & number of levels, $\lceil 2c^{*}\ln d\rceil$ & $33$ -- $59$ \\
$G$ & number of grid points on the $u$ axis & $3.6\cdot 10^4$ \\
$Q$ & quadrature order used to build tables (T1) & $96$ \\
$\mathcal{R}$ & set of count values for which tables are needed & $10^3$ -- $5\cdot 10^3$ values \\
$\mathcal{N}$ & number of states of a model & $10^3$ -- $2.5\cdot 10^5$ \\
$P$ & number of nonzero row entries, summed over rows & up to $1.5\cdot 10^7$ \\
$U$ & number of distinct profiles after deduplication & $2.6\cdot 10^4$ -- $7\cdot 10^4$ \\
$N$ & sum of one row's entries & $1$ -- $10^6$ \\
$k$ & number of nonzero entries of one row & $1$ -- $4{,}096$ \\
$\tilde k$ & number of distinct values in one row & $1$ -- $\sim 300$ \\
$C$ & number of table refreshes of a sequential code (T8) & $8$ -- $32$ \\
$D$ & number of lag predictors combined (T6, T7) & $6$ \\
$K$ & number of weight vectors in a combination grid & $\sim 10^2$ \\
\bottomrule
\end{tabular}
\end{center}

\section{T1 --- Building the tables of moment functions}

\field{Task}
Compute the moment functions $\varphi^{(L)}_r$ of
Eq.~\eqref{eq:phidef} numerically, for every level
$L \le L_{\max}$ and every count value $r$ that occurs in the data,
and store their logarithms on the $u$ grid.  This is done once per
alphabet size; the stored tables are then read by tasks T2 and T3.

\field{Formula}
Write $Y = Y' \cdot E_L$, where $Y'$ is the product of the first
$L - 1$ exponentials.  Taking the expectation over $E_L$ first gives
a recursion that connects level $\ell$ to level $\ell - 1$:
\begin{equation}
  \varphi^{(\ell)}_r(t)
  \;=\; \int_0^\infty e^{-v}\, v^{\,r}\,
  \varphi^{(\ell-1)}_r(t v)\, dv,
  \qquad
  \varphi^{(1)}_r(t) = \frac{\Gamma(r+1)}{(1+t)^{r+1}} .
  \label{eq:recursion}
\end{equation}

\field{Input}
The set $\mathcal{R}$ of count values that occur anywhere in the
data; the number of levels $L_{\max}$; the grid $u_1, \dots, u_G$.

\field{Output}
The numbers $\ln \varphi^{(L)}_r(e^{u_g})$ for all
$r \in \mathcal{R}$, $L \le L_{\max}$, $g \le G$ --- an array of
$|\mathcal{R}| \times L_{\max} \times G$ values, stored on disk.

\field{Algorithm}
For each $r$ separately: start from the closed form at level one;
then, for $\ell = 2, 3, \dots, L_{\max}$, evaluate the integral in
Eq.~\eqref{eq:recursion} at every grid point by Gauss--Laguerre
quadrature with $Q$ nodes.  The quadrature needs the previous level
at the shifted arguments $t v_j$; these values are obtained by linear
interpolation on the $u$ grid, using the exact limit
$\Gamma(r+1)^{\ell-1}$ to the left of the grid.  The work for
different values of $r$ is independent, so it is spread over all
processor cores; when the data later produces a count value that has
no table yet, that table is built on demand and added to the store.

\field{Cost and scaling}
Building: about $|\mathcal{R}| \cdot L_{\max} \cdot G \cdot Q$
interpolate-and-multiply steps; at the largest settings this is
$\approx 10^{12}$ steps, measured as $45$ minutes on $64$ cores.
Storage: $8 \cdot |\mathcal{R}| \cdot L_{\max} \cdot G$ bytes, which
reached $85$\,GB.  Growth: on a larger corpus, $|\mathcal{R}|$ grows
slowly (a new count value must appear for the first time, which
becomes rare); $G$ grows in proportion to the length of the $u$
interval, which grows like the logarithm of the largest count;
$L_{\max}$ grows like the logarithm of $d$.  The recurring cost after
building is \emph{reading} the tables, and the storage layout
determines whether reading is fast (a lesson learned on the shared
cluster file system).

\field{Ideas for speedup}
(1)~[exact] Keep one shared table store for all experiments; when we
audited the disk we found $219$\,GB of identical tables stored by
different runs.
(2)~[exact] Store the array level by level in single files, in
single precision: four times smaller, and a task that needs level
$L$ opens one file instead of $|\mathcal{R}|$ files.
(3)~[approx] For large $r$ the integrand of
Eq.~\eqref{eq:recursion} is sharply peaked and
$\varphi^{(L)}_r$ approaches a smooth limiting shape; beyond some
threshold the table could be replaced by this limiting formula.  The
error must be worked out before this is used.

\section{T2 --- Evaluating the probability of one profile}

\field{Task}
Given one profile, compute $\log q_L$ for every level $L$, and from
these $\log q_{\mathrm{avg}}$.  This is the innermost computation of
the project; we call it the \emph{atom}.  A static model invokes it
once per distinct profile; a sequential model invokes it again at
every table refresh (T8).

\field{Formula}
Eqs.~\eqref{eq:qdef}--\eqref{eq:psi}:
$\log q_L(\prof) = \log \int e^{\psi_L(u)}\,du \;-\; \ln \Gamma(N)$.

\field{Input}
The profile, given as its distinct values $c$ with multiplicities
$\mu_c$; the tables of T1.

\field{Output}
The vector $\bigl(\log_2 q_L(\prof)\bigr)_{L = 1..L_{\max}}$ and the
average $\log_2 q_{\mathrm{avg}}(\prof)$.

\field{Algorithm}
For each level: (i)~assemble the log-integrand $\psi_L$ on the grid
by Eq.~\eqref{eq:psi}; this reads $\tilde k + 1$ rows of the table
and costs $O(G \tilde k)$ additions.  (ii)~Find all local maxima of
$\psi_L$ that lie within a fixed window below the global maximum.
(The integrand can have two peaks: for deep levels and large counts a
second peak appears far to the left.)  (iii)~Refine each maximum by
solving $\psi_L'(u) = 0$ in a bracket around it, and add the peak
contributions by Laplace's method (a Gaussian approximation around
each peak, summed in the log domain).  (iv)~Add the closed-form left
tail from Section~\ref{sec:defs}.  Profiles with $N \le 1$, and the level $L = 1$,
are returned from closed forms without any of this.  Batches of
profiles are processed in parallel: the profiles are split over
worker processes, and each level's table rows are handed to the
workers through shared memory.

\field{Cost and scaling}
Step (i) dominates: $O(L_{\max}\, G\, \tilde k)$ per profile ---
between $4 \cdot 10^6$ operations for a poor profile ($\tilde k =
2$) and $2 \cdot 10^8$ for a rich one ($\tilde k \approx 100$).
Measured: $7 \cdot 10^4$ distinct profiles at $L_{\max} = 59$ in $17$
minutes on $64$ cores.

\field{Ideas for speedup}
(1)~[exact] \emph{Remember where the peaks are.}  The full sweep
over the grid exists only to locate the maxima reliably.  The
locations move very little when a profile changes slightly --- in
particular, between two refreshes of the same row.  Storing the peak
locations and starting the bracketed solve from the stored value
would let the algorithm evaluate $\psi_L$ only on a small window
around each peak, replacing $G$ by a window size of order $10^2$; a
full sweep remains as a fallback when the solve fails.  Potential
factor: up to $300$ on step (i).
(2)~[approx] \emph{Truncate the sum over levels.}  The sum in
Eq.~\eqref{eq:qdef} is dominated by a few levels around its largest
term, and the terms decay quickly away from it.  Evaluating the
dominant levels exactly and bounding the rest would replace
$L_{\max}$ by an effective count of $5$ to $10$.
(3)~[approx] Use coarser grids for profiles with small $N$, whose
integrands are wide and smooth; choose the grid per profile from a
target accuracy.

\section{T3 --- The prediction table of one row}

\field{Task}
Turn one row into a probability distribution for the symbol at the
next position.  This is the workhorse of the sequential codes of
T6--T8: at every refresh, every row of every predictor is converted
into such a table.

\field{Formula}
Let $\prof$ be the row's profile and let $c$ be the row's entry for
some symbol ($c = 0$ if the symbol is not observed in the row).
Write $\prof \oplus c$ for the profile obtained from $\prof$ by
increasing one entry with value $c$ by one (for $c = 0$: by
appending a new entry $1$).  The estimator's probability that the
next symbol is that particular symbol is
\begin{equation}
  p \;=\;
  \frac{q_{\mathrm{avg}}(\prof \oplus c)}{q_{\mathrm{avg}}(\prof)} .
  \label{eq:pred}
\end{equation}
The value depends on the symbol only through its count $c$, so the
table has only $\tilde k + 1$ distinct entries (one per distinct
count value, and one shared by all unobserved symbols).  The entries
sum to one exactly; in practice they are renormalized to absorb the
numerical error of T2, measured at about $5 \cdot 10^{-5}$ relative.

\field{Input}
The row's counts; the tables of T1.

\field{Output}
The logarithms of the $\tilde k + 1$ distinct probabilities, and the
assignment of each symbol to its probability.

\field{Algorithm}
Today: one T2 evaluation for $\prof$ and one for each of the
$\tilde k + 1$ profiles $\prof \oplus c$; results are cached by
profile, so a profile that appears again (in another row, or at
another refresh) is not recomputed.

\field{Cost and scaling}
$\tilde k + 2$ atom evaluations per row, i.e.\ $O(L_{\max} G \tilde
k^2)$ --- note the square.  Measured: the pooled run at $d =
1{,}024$ with $C = 32$ refreshes needed $5.9 \cdot 10^6$ atom
evaluations, which is the dominant cost of that run; the cache helps
little because an active row's profile is different at every
refresh.

\field{Ideas for speedup}
(1)~[exact; the largest single lever in this document]
\emph{Compute all the ratios from one integrand.}  The profiles
$\prof \oplus c$ differ from $\prof$ in a single factor of the
integrand in Eq.~\eqref{eq:integral}.  Concretely,
\[
  q_L(\prof \oplus c) \;=\; \frac{1}{\Gamma(N+1)} \int
  e^{\psi_L(u)}\, \cdot\, e^{u}\,
  \frac{\varphi^{(L)}_{c+1}(e^u)}{\varphi^{(L)}_{c}(e^u)}\; du ,
\]
with $\psi_L$ the log-integrand of the \emph{base} profile,
Eq.~\eqref{eq:psi}.  So: assemble $\psi_L$ once, at cost
$O(G\tilde k)$; then every one of the $\tilde k + 1$ entries is a
single weighted sum over the grid, at cost $O(G)$ each.  The cost
per row falls from $O(G \tilde k^{2})$ to $O(G \tilde k)$ --- a
factor $\tilde k$, which is about $100$ exactly for the rich rows
that dominate --- and the result is the same numbers.
(2)~[exact] \emph{Update the integrand instead of rebuilding it.}
Between two refreshes only a few entries of a row change.  If the
row's $\psi_L$ is kept in memory, it can be updated by adding
$\ln \varphi^{(L)}_{c'} - \ln \varphi^{(L)}_{c}$ for each changed
entry: cost proportional to the number of changes, not to $\tilde
k$.  Keeping $\psi_L$ for one row costs $G \times L_{\max}$ single-
precision numbers ($\approx 8$\,MB at the largest settings), which
is affordable for the busy rows --- and the busy rows are exactly
the expensive ones.
(3)~The ideas of T2 apply on top.

\section{T4 --- Counting}

\field{Task}
Read the corpus once and produce, for a given state rule $\sigma$,
all rows, their profiles, and the assignment of each row to its
profile.

\field{Formula}
$m_x(s)$ as in Section~\ref{sec:defs}, for every state $s$ that occurs; then
profiles and deduplication.

\field{Input}
The corpus as $n$ integers; the rule $\sigma$.

\field{Output}
$\mathcal{N}$ rows with $P$ nonzero entries in total; the list of
$U$ distinct profiles; the map from rows to profiles.

\field{Algorithm}
One pass over the corpus with a hash table from (state, symbol)
pairs to counters; then one pass over rows to sort each row's
nonzero entries and to deduplicate profiles.

\field{Cost and scaling}
Time is linear in $n$ and negligible (minutes).  \emph{Memory} is
the problem: the hash table costs $30$ to $60$ bytes per nonzero
entry, measured at about $200$\,GB for $P \approx 1.5\cdot 10^7$
(all contexts of depth up to two at full alphabet).  This is what
currently forbids depth three.  After deduplication the downstream
tasks see only $U \ll P$ profiles; measured examples: $4.4 \cdot
10^6$ rows collapse to $7 \cdot 10^4$ profiles.

\field{Ideas for speedup}
(1)~[exact] Count on disk: sort the (state, symbol) pairs
externally and merge equal pairs.  Standard technology; memory
becomes constant, time $O(n \log n)$ with small constants.
(2)~[exact] Build a suffix array of the corpus (a sorted index of
all its positions, storable in a few bytes per position).  Every
context of every depth corresponds to one contiguous interval of
the suffix array, so the counts for \emph{all} depths at once can be
read off from it.  This is the natural route to depth three and
beyond.
(3)~[exact] Replace hash tables of Python objects by fixed-width
integer arrays; this alone reduces the memory constant by an order
of magnitude.

\section{T5 --- Aggregation}

\field{Task}
Combine the atoms into the reported numbers: the compressed length
of a model, the mixture over a family of models, the posterior over
the family.

\field{Formula}
The length of one model is $\sum_{\text{rows}} -\log_2
q_{\mathrm{avg}}(\text{profile of the row})$.  For the context-tree
family the mixture over all prunings is computed by the recursion
$\beta(v) = \frac12 q_{\mathrm{avg}}(\prof_v) + \frac12 \prod_{w}
\beta(w)$, where $v$ runs over observed contexts, $\prof_v$ is the
row of $v$, and $w$ runs over the observed one-symbol extensions of
$v$.  The mixture over any finite family of $K$ models is the
log-sum-exp of their total log-probabilities minus $\log_2 K$, and
the posterior over the family is the normalized exponential of the
same totals.

\field{Input/Output/Algorithm/Cost}
Inputs are the per-profile atoms (T2) and the context structure;
outputs are bits-per-symbol figures and posteriors.  Single passes
over rows and contexts; cost $O(\mathcal{N})$ or $O(P)$; measured in
seconds.  No scaling concern.

\section{T6 --- Combining predictors by weighted sums}

\field{Task}
Given $D$ lag predictors and the no-memory predictor, code the
corpus under \emph{every} candidate weight vector of a grid
simultaneously, in one pass.

\field{Formula}
At position $t$, predictor $e$ supplies the probability
$p_{e,t}$ that it assigned to the symbol that actually occurred at
$t$ (read from its current prediction table, T3).  For each weight
vector $\lambda = (\lambda_0, \dots, \lambda_D)$ of a grid of size
$K$ (nonnegative entries summing to one):
\[
  \mathrm{bits}(\lambda) \;=\; \sum_{t}\; -\log_2 \Bigl(\,
  \sum_{e=0}^{D} \lambda_e\; p_{e,t} \Bigr).
\]

\field{Input}
The $n \times (D+1)$ array of the numbers $p_{e,t}$; the grid of
weight vectors.

\field{Output}
$K$ compressed lengths, one per weight vector; their mixture and
posterior (T5).

\field{Algorithm and cost}
One matrix product of the $p$ array with the weight matrix per block
of positions, then logarithms: $O(n (D+1) K)$ multiplications,
about $10^{10}$ at our settings --- minutes.  Linear in everything;
never a bottleneck.

\section{T7 --- Combining predictors by tempered products}

\field{Task}
The same as T6 for the multiplicative rule.  The product of
probabilities must be renormalized \emph{over the whole alphabet at
every position}, and this normalization is the expensive part.

\field{Formula}
Let $q_0$ be the no-memory predictor's table and $p_j(x)$ the
probability that lag predictor $j$ currently assigns to symbol $x$.
For an exponent vector $\beta = (\beta_1, \dots, \beta_D)$ from a
grid of size $K_{\mathrm{prod}}$:
\[
  p_\beta(x) \;=\; \frac{s_\beta(x)}{Z}, \qquad
  s_\beta(x) \;=\; q_0(x) \prod_{j=1}^{D}
  \biggl(\frac{p_j(x)}{q_0(x)}\biggr)^{\!\beta_j}, \qquad
  Z \;=\; \sum_{x=1}^{d} s_\beta(x),
\]
where $Z$ is recomputed at every position (the tables $p_j$ depend
on the current states).

\field{Input}
The prediction tables (T3) of the $D$ lag predictors and of $q_0$;
the exponent grid.

\field{Output}
$K_{\mathrm{prod}}$ compressed lengths.

\field{Algorithm today, and cost}
Dense: at every position, form the $d$ numbers $\ln s_\beta(x)$ and
sum their exponentials.  Cost $O(n\, d\, D\, K_{\mathrm{prod}})$:
measured $68$ minutes at $d = 4{,}096$; at $d = 10^6$ and $n =
1.5\cdot 10^8$ it would be $\sim 10^{16}$ operations.  \textbf{This
is the only part of the project whose cost is proportional to the
alphabet size at every position.}

\field{Idea for speedup}
[exact; removes the factor $d$]  By Eq.~\eqref{eq:pred}, each
prediction table consists of one shared value for all symbols not
observed in the row, plus individual values for the observed ones.
Let $S$ be the union, over the $D$ predictors, of the sets of
symbols observed in their current rows; $|S|$ is a few hundred to a
few thousand, far below $d$.  For every symbol outside $S$, each
factor of $s_\beta(x)$ takes its shared value, so $s_\beta(x) =
(\text{const}) \cdot q_0(x)^{\gamma}$ with $\gamma = 1 - \sum_j
\beta_j$.  Therefore
\[
  Z \;=\; (\text{const}) \cdot \Bigl( T_\gamma -
  \sum_{x \in S} q_0(x)^{\gamma} \Bigr)
  \;+\; \sum_{x \in S} s_\beta(x),
  \qquad
  T_\gamma \;=\; \sum_{x=1}^{d} q_0(x)^{\gamma},
\]
and the power sums $T_\gamma$ are computed once per refresh.  The
normalization then costs $O(|S| D)$ per position instead of
$O(d\,D)$ --- exact, and it makes the multiplicative rule feasible
at any alphabet size.

\section{T8 --- Scheduling the refreshes of a sequential code}

\field{Task}
A sequential code reads the corpus once and predicts each position
from the counts gathered so far.  Recomputing every prediction table
at every position would be exact but hopeless; instead the tables
are \emph{refreshed} at chosen moments and used unchanged in
between.  Any refresh schedule that depends only on the past yields
a valid code; using slightly outdated tables only lengthens the
output a little.  The schedule determines how many times T3 runs,
and is therefore the multiplier that separates static runs
(minutes) from sequential runs (hours).

\field{Current schedule and cost}
A global clock: all rows are refreshed at $C$ evenly spaced
positions.  T3 work: $C$ times the number of active rows; measured
$5.9\cdot 10^6$ atom evaluations at $d = 1{,}024$, $C = 32$.

\field{Ideas for speedup}
(1)~[valid code; staleness controlled] \emph{Refresh on growth.}
Rebuild a row's table when the row's total $N$ has grown by a fixed
factor $1 + \epsilon$ since the last rebuild.  A row that ends with
total $N$ is then rebuilt only $\log_{1+\epsilon} N$ times ---
logarithmic instead of $C$ --- and the table a position actually
uses is never based on counts more than a factor $1 + \epsilon$
behind, which is the kind of error that matters for prediction
quality.  The trigger depends only on the past, so the code remains
valid.  Combined with ideas T3(1) and T3(2), the estimated overall
gain on sequential runs is two to three orders of magnitude.
(2)~[exact] Distribute rows over cluster nodes; rows are independent
once the schedule is fixed.

\section*{Scaling summary}

Consider a corpus ten times larger ($n \approx 1.4\cdot 10^8$,
alphabet $\sim 10^6$).  With today's algorithms: T4 would need
terabytes of memory (it already blocks depth three); T3 under the
schedule of T8 would take weeks; T7 in its dense form is out of the
question; T1 and T2 grow only through $|\mathcal{R}|$, $G$ and
$L_{\max}$ and remain manageable.  With the ideas marked [exact] ---
one shared integrand per row (T3), integrand updates (T3), remembered
peaks (T2), the sparse normalization (T7), refresh on growth (T8),
counting on disk or via a suffix array (T4), and distribution over
nodes --- each task returns to the scale of node-hours.  The ideas
marked [approx] (level truncation, limiting formulas for large
counts, coarser grids) are held in reserve, each with an error that
can be checked against the exact computation on samples.

Suggested order of discussion: T3(1), T7, T8(1), T2(1), T4(2); the
approximations only if something is still too slow after these.

\end{document}
